% Options for packages loaded elsewhere
\PassOptionsToPackage{unicode}{hyperref}
\PassOptionsToPackage{hyphens}{url}
\documentclass[
]{article}
\usepackage{xcolor}
\usepackage{amsmath,amssymb}
\setcounter{secnumdepth}{-\maxdimen} % remove section numbering
\usepackage{iftex}
\ifPDFTeX
  \usepackage[T1]{fontenc}
  \usepackage[utf8]{inputenc}
  \usepackage{textcomp} % provide euro and other symbols
\else % if luatex or xetex
  \usepackage{unicode-math} % this also loads fontspec
  \defaultfontfeatures{Scale=MatchLowercase}
  \defaultfontfeatures[\rmfamily]{Ligatures=TeX,Scale=1}
\fi
\usepackage{lmodern}
\ifPDFTeX\else
  % xetex/luatex font selection
\fi
% Use upquote if available, for straight quotes in verbatim environments
\IfFileExists{upquote.sty}{\usepackage{upquote}}{}
\IfFileExists{microtype.sty}{% use microtype if available
  \usepackage[]{microtype}
  \UseMicrotypeSet[protrusion]{basicmath} % disable protrusion for tt fonts
}{}
\makeatletter
\@ifundefined{KOMAClassName}{% if non-KOMA class
  \IfFileExists{parskip.sty}{%
    \usepackage{parskip}
  }{% else
    \setlength{\parindent}{0pt}
    \setlength{\parskip}{6pt plus 2pt minus 1pt}}
}{% if KOMA class
  \KOMAoptions{parskip=half}}
\makeatother
\usepackage{longtable,booktabs,array}
\newcounter{none} % for unnumbered tables
\usepackage{calc} % for calculating minipage widths
% Correct order of tables after \paragraph or \subparagraph
\usepackage{etoolbox}
\makeatletter
\patchcmd\longtable{\par}{\if@noskipsec\mbox{}\fi\par}{}{}
\makeatother
% Allow footnotes in longtable head/foot
\IfFileExists{footnotehyper.sty}{\usepackage{footnotehyper}}{\usepackage{footnote}}
\makesavenoteenv{longtable}
\setlength{\emergencystretch}{3em} % prevent overfull lines
\providecommand{\tightlist}{%
  \setlength{\itemsep}{0pt}\setlength{\parskip}{0pt}}
\usepackage{bookmark}
\IfFileExists{xurl.sty}{\usepackage{xurl}}{} % add URL line breaks if available
\urlstyle{same}
% fallback for those not using the hyperref driver hyperxmp:
\makeatletter
\@ifundefined{xmpquote}{\newcommand{\xmpquote}[1]{#1}}{}
\makeatother
\hypersetup{
  hidelinks,
  pdfcreator={LaTeX via pandoc}}

\author{}
\date{}

\begin{document}

\section{UNIVERSIDADE DE BRASÍLIA}\label{universidade-de-brasuxedlia}

\subsection{Faculdade de Administração, Contabilidade, Economia e Gestão
de Políticas
Públicas}\label{faculdade-de-administrauxe7uxe3o-contabilidade-economia-e-gestuxe3o-de-poluxedticas-puxfablicas}

\section{RELATÓRIO TÉCNICO DE AUDITORIA
CONTÁBIL}\label{relatuxf3rio-tuxe9cnico-de-auditoria-contuxe1bil}

\textbf{Estudo acadêmico --- Auditoria 1}

\textbf{Entidade examinada:} Equatorial Energia Distribuição S.A.

\textbf{Exercício social:} 2025

\textbf{Área contábil:} Ativo financeiro da concessão e obrigações
especiais relacionadas

\begin{center}\rule{0.5\linewidth}{0.5pt}\end{center}

\textbf{Aluno:} Cláudio Wollace Santos Da Hora\\
\textbf{Matrícula:} 242010937\\
\textbf{Professor:} Lucas Oliveira Gomes Ferreira\\
\textbf{Disciplina:} Auditoria 1

\textbf{Data de elaboração:} 25 de outubro de 2026

\textbf{Natureza do trabalho:} Estudo acadêmico de auditoria
independente, com análise documental pública e procedimentos analíticos
sobre base de dados sintética.

\textbf{Situação:} Versão preliminar, sujeita à verificação das fontes,
revisão das normas aplicáveis e conclusão dos papéis de trabalho.

\begin{center}\rule{0.5\linewidth}{0.5pt}\end{center}

\section{1. Identificação do
trabalho}\label{identificauxe7uxe3o-do-trabalho}

Este relatório apresenta um estudo acadêmico de auditoria contábil sobre
o ativo financeiro da concessão da Equatorial Energia Distribuição S.A.,
com referência às demonstrações contábeis do exercício findo em 31 de
dezembro de 2025.

O trabalho foi desenvolvido no âmbito da disciplina Auditoria 1 e tem
como finalidade aplicar conceitos e procedimentos relacionados ao
planejamento de auditoria, à identificação e avaliação de riscos, à
materialidade, à obtenção e avaliação de evidências, à amostragem e à
documentação dos resultados.

A análise utiliza informações públicas da entidade e uma base de dados
sintética elaborada especificamente para o exercício acadêmico. Os
resultados obtidos com essa base representam testes de metodologia e não
constituem constatações sobre os registros reais da companhia.

\section{2. Objetivo, escopo e natureza do
trabalho}\label{objetivo-escopo-e-natureza-do-trabalho}

\subsection{2.1. Objetivo geral}\label{objetivo-geral}

O objetivo geral é desenvolver um estudo técnico de auditoria sobre o
ativo financeiro da concessão, relacionando os saldos divulgados pela
entidade aos riscos de distorção relevante, às afirmações contábeis
pertinentes e aos procedimentos que poderiam ser planejados para obter
evidência de auditoria apropriada e suficiente.

O trabalho também busca demonstrar a aplicação acadêmica das Normas
Brasileiras de Contabilidade relacionadas à auditoria independente,
abrangendo planejamento, avaliação de riscos, materialidade, evidência,
amostragem, documentação e avaliação das distorções identificadas.

\subsection{2.2. Objetivos específicos}\label{objetivos-especuxedficos}

\begin{itemize}
\tightlist
\item
  Identificar informações públicas relevantes sobre a conta contábil
  selecionada.
\item
  Examinar os saldos divulgados do ativo financeiro da concessão e das
  obrigações especiais relacionadas.
\item
  Identificar riscos preliminares de auditoria e afirmações contábeis
  relevantes.
\item
  Definir parâmetros acadêmicos de materialidade e documentar suas
  premissas.
\item
  Desenvolver uma base sintética com pelo menos 500.000 registros
  analíticos únicos, conciliada aos saldos públicos de referência.
\item
  Executar testes automatizados sobre a base sintética.
\item
  Documentar critérios, resultados, exceções e limitações dos testes.
\item
  Comparar as inconsistências deliberadamente inseridas com as exceções
  detectadas.
\item
  Elaborar papéis de trabalho, programas de auditoria e matriz de
  cobertura das normas aplicáveis.
\item
  Distinguir análises documentais públicas, procedimentos executados
  sobre dados sintéticos e procedimentos apenas planejados.
\end{itemize}

\subsection{2.3. Escopo}\label{escopo}

O escopo principal concentra-se na informação consolidada apresentada na
Nota Explicativa nº 11, referente ao ativo financeiro da concessão,
conforme o documento de demonstrações contábeis de 2025 utilizado como
referência.

O trabalho contempla:

\begin{enumerate}
\def\labelenumi{\arabic{enumi}.}
\tightlist
\item
  Identificação dos saldos públicos de referência.
\item
  Caracterização da conta contábil e dos riscos associados.
\item
  Conciliação dos valores de referência com a base sintética.
\item
  Testes automatizados de integridade, recálculo, documentação de
  suporte e período de referência.
\item
  Avaliação ilustrativa de amostragem e materialidade.
\item
  Documentação das exceções simuladas e dos resultados dos testes.
\item
  Planejamento de procedimentos complementares que exigiriam acesso aos
  registros e documentos originais da entidade.
\end{enumerate}

Este trabalho não constitui uma auditoria efetivamente realizada na
companhia, não substitui o relatório de auditor independente e não
fornece opinião de auditoria sobre as demonstrações contábeis da
Equatorial Energia Distribuição S.A.

\subsection{2.4. Classificação dos
procedimentos}\label{classificauxe7uxe3o-dos-procedimentos}

Para assegurar transparência metodológica, os procedimentos são
classificados em três categorias.

\textbf{a) Análise documental pública:} leitura e extração de
informações constantes das demonstrações contábeis e de outras fontes
públicas identificadas.

\textbf{b) Procedimentos analíticos sobre dados sintéticos:} execução de
testes em base artificial criada para o trabalho acadêmico, com regras
de geração documentadas e inconsistências deliberadamente inseridas para
avaliar a capacidade de detecção dos testes.

\textbf{c) Procedimentos planejados, não executados em registros reais:}
atividades que exigiriam acesso ao razão contábil, documentos originais,
sistemas internos, confirmações externas ou responsáveis pela entidade.
Esses procedimentos são apresentados como planejamento, sem alegação de
execução efetiva.

Os resultados obtidos na base sintética não devem ser interpretados como
evidência de erros reais, deficiências de controle ou irregularidades na
companhia.

\section{3. Fontes de informação, critérios e
metodologia}\label{fontes-de-informauxe7uxe3o-crituxe9rios-e-metodologia}

\subsection{3.1. Fonte documental
principal}\label{fonte-documental-principal}

A fonte documental principal é o arquivo de demonstrações contábeis
anuais e notas explicativas referente ao exercício de 2025,
disponibilizado para a elaboração deste trabalho.

\textbf{Referência utilizada:} Nota Explicativa nº 11 --- Ativo
financeiro da concessão, página impressa 44, correspondente à página 46
do arquivo PDF, conforme o controle de fontes do projeto.

Os valores extraídos dessa nota são utilizados como referências públicas
para a construção e a conciliação da base sintética. Antes da
finalização do relatório, deverão ser confirmados o endereço eletrônico
de origem, a identificação completa do documento e sua correspondência
com a versão publicada pela companhia.

O controle de fontes deverá registrar, para cada documento, título,
entidade responsável, exercício, endereço eletrônico, data de acesso,
página impressa, página do arquivo e finalidade da utilização.

\subsection{3.2. Normas e critérios de
auditoria}\label{normas-e-crituxe9rios-de-auditoria}

O trabalho utiliza como referencial as Normas Brasileiras de
Contabilidade aplicáveis à auditoria independente, incluindo a Estrutura
Conceitual e as NBC TA 200, 210, 220, 230, 240, 250, 260, 265, 300, 315,
320, 330, 402, 450, 500, 501, 505, 510, 520, 530, 540, 550, 560, 570,
580 e 600, conforme o roteiro acadêmico.

A versão vigente e a aplicabilidade de cada norma deverão ser
verificadas na fonte oficial do Conselho Federal de Contabilidade. A
matriz de cobertura deverá indicar a norma, os itens pertinentes, o
procedimento associado, a evidência correspondente e a situação de
atendimento.

A inclusão de uma norma na matriz não significa que todos os seus
requisitos tenham sido executados ou integralmente atendidos. A
documentação deverá distinguir entre análise conceitual, procedimento
planejado, procedimento executado em dados sintéticos e requisito
pendente de verificação.

\subsection{3.3. Construção da base
sintética}\label{construuxe7uxe3o-da-base-sintuxe9tica}

Para possibilitar testes quantitativos reproduzíveis sem acesso aos
registros internos da companhia, foi desenvolvida uma base de dados
sintética com 500.000 registros únicos.

A base foi construída para reproduzir, em nível agregado, os saldos
públicos de referência da conta examinada. Os registros são artificiais
e foram gerados por código, com identificação explícita de sua natureza
sintética.

A estrutura contempla componentes do ativo financeiro da concessão e das
obrigações especiais relacionadas, datas, categorias de registros,
valores registrados, valores recalculados e campos de documentação de
suporte. A base inclui inconsistências deliberadamente inseridas para
testar os procedimentos automatizados.

A descrição das variáveis, os critérios de geração, as premissas, os
códigos de execução e os controles de integridade devem ser mantidos nos
arquivos de dados, dicionário de variáveis, scripts e papéis de trabalho
do projeto.

\textbf{Advertência metodológica:} a base sintética não corresponde ao
razão contábil da companhia, não contém transações reais identificadas e
não permite concluir que os registros, documentos ou exceções simulados
existam nos sistemas da entidade.

\subsection{3.4. Procedimentos de
análise}\label{procedimentos-de-anuxe1lise}

A metodologia combina análise documental, conciliação de saldos de
referência, testes automatizados e documentação estruturada dos
resultados.

Os procedimentos sobre a base sintética incluem:

\begin{itemize}
\tightlist
\item
  Verificação da quantidade de registros e da unicidade dos
  identificadores.
\item
  Conferência da classificação dos registros como sintéticos.
\item
  Conciliação dos valores agregados aos saldos de referência definidos
  para o exercício.
\item
  Comparação entre valores registrados e valores recalculados.
\item
  Identificação de registros sem indicação de documento de suporte.
\item
  Identificação de datas fora do período de referência.
\item
  Comparação entre as inconsistências deliberadamente inseridas e as
  exceções detectadas.
\item
  Registro dos resultados, limitações e referências cruzadas nos papéis
  de trabalho.
\end{itemize}

A amostragem ilustrativa é utilizada como exercício metodológico. Sua
seleção e seus resultados não devem ser confundidos com uma amostra de
registros reais da companhia nem com evidência suficiente para sustentar
uma opinião de auditoria sobre as demonstrações contábeis.

\subsection{3.5. Reprodutibilidade e
documentação}\label{reprodutibilidade-e-documentauxe7uxe3o}

Os procedimentos foram organizados em scripts de R, arquivos de
resultados, papéis de trabalho e controles de fontes e normas. A
documentação deve permitir que outro leitor compreenda as premissas
utilizadas, reproduza os testes e identifique a origem de cada valor
apresentado.

Os resultados somente serão considerados concluídos após a conferência
dos arquivos gerados, a revisão das fórmulas e dos critérios de teste e
a validação das referências cruzadas. A revisão independente permanece
pendente quando não houver evidência documental de sua realização.

\section{4. Caracterização da conta contábil
examinada}\label{caracterizauxe7uxe3o-da-conta-contuxe1bil-examinada}

\subsection{4.1. Natureza da conta}\label{natureza-da-conta}

O ativo financeiro da concessão está relacionado aos direitos econômicos
associados aos contratos de concessão e ao tratamento contábil aplicável
aos investimentos e à remuneração vinculados à prestação de serviços
públicos concedidos.

A avaliação dessa conta exige compreender os termos contratuais, as
regras regulatórias aplicáveis, os critérios de reconhecimento e
mensuração, a movimentação dos saldos e as divulgações efetuadas nas
demonstrações contábeis.

A análise de auditoria deve considerar se os valores reconhecidos estão
adequadamente fundamentados, se as movimentações foram registradas no
período correto e se os critérios contábeis e as estimativas utilizadas
são apropriados. A natureza exata dos direitos reconhecidos e dos
critérios de mensuração deverá ser confirmada nas políticas contábeis e
nas notas explicativas completas da entidade.

\subsection{4.2. Saldos públicos de
referência}\label{saldos-puxfablicos-de-referuxeancia}

Conforme os valores identificados na Nota Explicativa nº 11, os saldos
consolidados utilizados como referência para o exercício de 2025 são:

{\def\LTcaptype{none} % do not increment counter
\begin{longtable}[]{@{}lr@{}}
\toprule\noalign{}
Componente & Valor (R\$ mil) \\
\midrule\noalign{}
\endhead
\bottomrule\noalign{}
\endlastfoot
Ativo financeiro da concessão --- valor bruto & 21.062.007 \\
Obrigações especiais & (5.862.808) \\
\textbf{Ativo financeiro da concessão --- valor líquido} &
\textbf{15.199.199} \\
\end{longtable}
}

\textbf{Fonte:} demonstrações contábeis de 2025, Nota Explicativa nº 11,
página impressa 44, página 46 do arquivo PDF.

Os valores estão expressos em milhares de reais, conforme a apresentação
do documento de referência. O saldo líquido resulta da dedução das
obrigações especiais do valor bruto do ativo financeiro da concessão.

A tabela deverá ser conferida novamente com o documento original antes
da emissão da versão final do relatório. Os valores acima representam
referências documentais públicas e não constituem, isoladamente,
evidência de auditoria sobre a existência, a exatidão ou a avaliação dos
saldos.

\subsection{4.3. Relevância para a
auditoria}\label{relevuxe2ncia-para-a-auditoria}

O saldo líquido de R\$ 15.199.199 mil é expressivo e justifica a
priorização acadêmica da conta para o desenvolvimento dos procedimentos
analíticos. A materialidade, entretanto, deve ser definida e documentada
conforme o referencial adotado, considerando também fatores
qualitativos, riscos específicos e as necessidades dos usuários das
demonstrações contábeis.

A magnitude do saldo não demonstra, por si só, a existência de
distorção. Ela indica que eventuais erros relevantes na classificação,
mensuração, reconhecimento ou divulgação poderiam afetar a interpretação
das informações contábeis.

\subsection{4.4. Afirmações contábeis
relevantes}\label{afirmauxe7uxf5es-contuxe1beis-relevantes}

Para o planejamento dos procedimentos, são consideradas as seguintes
afirmações:

\begin{itemize}
\tightlist
\item
  \textbf{Existência:} verificar, em uma auditoria real, se os direitos
  registrados correspondem a ativos existentes e sustentados por
  contratos, informações regulatórias e documentação pertinente.
\item
  \textbf{Direitos e obrigações:} avaliar se a entidade detém os
  direitos reconhecidos e se as obrigações especiais foram identificadas
  e apresentadas adequadamente.
\item
  \textbf{Integralidade:} avaliar se os ativos, as obrigações e as
  movimentações relevantes foram registrados de forma completa.
\item
  \textbf{Exatidão, avaliação e alocação:} examinar a correção dos
  cálculos, a aplicação dos critérios contábeis e a razoabilidade das
  estimativas utilizadas.
\item
  \textbf{Corte:} verificar se adições, baixas, transferências e demais
  movimentações foram reconhecidas no período apropriado.
\item
  \textbf{Classificação e apresentação:} avaliar se os saldos estão
  classificados e apresentados adequadamente nas demonstrações
  contábeis.
\item
  \textbf{Divulgação:} verificar se as notas explicativas contêm
  informações suficientes e consistentes com os critérios contábeis
  aplicáveis.
\end{itemize}

Essas afirmações orientam o planejamento, mas não são consideradas
comprovadas apenas pela conciliação matemática da base sintética.

\subsection{4.5. Principais riscos
preliminares}\label{principais-riscos-preliminares}

Os riscos abaixo são hipóteses de planejamento, não constatações de
auditoria:

\begin{enumerate}
\def\labelenumi{\arabic{enumi}.}
\tightlist
\item
  \textbf{Mensuração inadequada:} aplicação incorreta de premissas ou
  critérios contábeis na determinação dos saldos.
\item
  \textbf{Inconsistência nos cálculos:} divergências entre valores
  registrados e valores recalculados.
\item
  \textbf{Registro incompleto:} omissão de adições, baixas, obrigações
  especiais ou outras movimentações relevantes.
\item
  \textbf{Erro de corte:} reconhecimento de movimentações em período
  diferente daquele a que pertencem.
\item
  \textbf{Classificação inadequada:} apresentação incorreta de
  componentes do ativo ou das obrigações relacionadas.
\item
  \textbf{Suporte documental insuficiente:} ausência ou inadequação de
  evidências que sustentem os registros.
\item
  \textbf{Divulgação insuficiente:} falta de informações relevantes
  sobre critérios contábeis, premissas ou movimentações.
\end{enumerate}

A avaliação desses riscos deverá ser aprofundada com base nas notas
explicativas, nas políticas contábeis, nos contratos e nas demais
informações públicas disponíveis. Procedimentos que dependam de
registros internos ou documentação original permanecerão classificados
como planejados, caso não possam ser executados.

\section{5. Materialidade e avaliação preliminar de
riscos}\label{materialidade-e-avaliauxe7uxe3o-preliminar-de-riscos}

\subsection{5.1. Finalidade da
materialidade}\label{finalidade-da-materialidade}

A materialidade é um parâmetro utilizado no planejamento e na execução
da auditoria para orientar a natureza, a época e a extensão dos
procedimentos, bem como a avaliação das distorções identificadas
individualmente e em conjunto.

A definição de materialidade em uma auditoria real exige julgamento
profissional e consideração das necessidades dos usuários das
demonstrações contábeis, da natureza da entidade, de seus resultados e
de outros fatores relevantes. Como este trabalho é acadêmico e não
dispõe de todas as informações necessárias para uma avaliação
profissional completa, os parâmetros apresentados a seguir são
exclusivamente ilustrativos.

\subsection{5.2. Materialidade global
ilustrativa}\label{materialidade-global-ilustrativa}

Para fins didáticos, adotou-se como referência o saldo líquido do ativo
financeiro da concessão de R\$ 15.199.199 mil, conforme os valores
públicos utilizados neste trabalho.

Foi aplicada uma taxa acadêmica de 1\% sobre esse saldo:

{\def\LTcaptype{none} % do not increment counter
\begin{longtable}[]{@{}
  >{\raggedright\arraybackslash}p{(\linewidth - 4\tabcolsep) * \real{0.3000}}
  >{\raggedright\arraybackslash}p{(\linewidth - 4\tabcolsep) * \real{0.3000}}
  >{\raggedleft\arraybackslash}p{(\linewidth - 4\tabcolsep) * \real{0.4000}}@{}}
\toprule\noalign{}
\begin{minipage}[b]{\linewidth}\raggedright
Parâmetro
\end{minipage} & \begin{minipage}[b]{\linewidth}\raggedright
Critério
\end{minipage} & \begin{minipage}[b]{\linewidth}\raggedleft
Valor (R\$ mil)
\end{minipage} \\
\midrule\noalign{}
\endhead
\bottomrule\noalign{}
\endlastfoot
Base de referência & Saldo líquido do ativo financeiro da concessão &
15.199.199,00 \\
Materialidade global ilustrativa & 1\% da base & 151.991,99 \\
Materialidade de execução ilustrativa & 75\% da materialidade global &
113.993,99 \\
Limite trivial ilustrativo & 5\% da materialidade global & 7.599,60 \\
\end{longtable}
}

Os percentuais foram escolhidos para fins acadêmicos e não representam
parâmetros oficiais da companhia, de seus auditores ou de qualquer órgão
regulador.

A utilização do saldo de uma conta patrimonial como base de referência
também é uma escolha metodológica do exercício. Em um trabalho real,
seria necessário avaliar se essa base é apropriada diante das
circunstâncias e dos usuários das demonstrações contábeis.

\subsection{5.3. Interpretação dos
parâmetros}\label{interpretauxe7uxe3o-dos-paruxe2metros}

A materialidade global ilustrativa serve como referência para avaliar a
possível relevância quantitativa das distorções. A materialidade de
execução foi estabelecida em 75\% da materialidade global para
demonstrar como o auditor pode adotar um limite inferior no planejamento
dos procedimentos.

O limite trivial foi calculado em 5\% da materialidade global para fins
de organização e avaliação didática das diferenças. Esses parâmetros não
dispensam a análise qualitativa: uma distorção inferior aos limites
quantitativos pode ser relevante em razão de sua natureza ou das
circunstâncias em que ocorreu.

Os valores devem ser reavaliados se a base de referência, as premissas
ou o escopo do trabalho forem alterados.

\subsection{5.4. Matriz preliminar de
riscos}\label{matriz-preliminar-de-riscos}

{\def\LTcaptype{none} % do not increment counter
\begin{longtable}[]{@{}
  >{\raggedright\arraybackslash}p{(\linewidth - 4\tabcolsep) * \real{0.3333}}
  >{\raggedright\arraybackslash}p{(\linewidth - 4\tabcolsep) * \real{0.3333}}
  >{\raggedright\arraybackslash}p{(\linewidth - 4\tabcolsep) * \real{0.3333}}@{}}
\toprule\noalign{}
\begin{minipage}[b]{\linewidth}\raggedright
Risco
\end{minipage} & \begin{minipage}[b]{\linewidth}\raggedright
Afirmações relacionadas
\end{minipage} & \begin{minipage}[b]{\linewidth}\raggedright
Possível resposta de auditoria
\end{minipage} \\
\midrule\noalign{}
\endhead
\bottomrule\noalign{}
\endlastfoot
Erro de mensuração ou cálculo & Exatidão, avaliação e alocação &
Recalcular valores e comparar os resultados com a documentação de
origem \\
Omissão de movimentações & Integralidade & Examinar a movimentação da
conta e conciliar os saldos com fontes independentes \\
Registro em período incorreto & Corte & Inspecionar, em auditoria real,
documentos próximos ao encerramento do exercício \\
Classificação inadequada & Classificação e apresentação & Comparar a
classificação contábil com a política aplicável e a documentação
pertinente \\
Falta de suporte documental & Existência, exatidão e avaliação &
Inspecionar contratos, memórias de cálculo e documentos originais \\
Estimativas inadequadas & Avaliação e alocação & Avaliar métodos,
premissas, dados utilizados e incertezas relevantes \\
Divulgação insuficiente & Apresentação e divulgação & Confrontar as
notas explicativas com os requisitos contábeis aplicáveis \\
\end{longtable}
}

Esta matriz apresenta riscos preliminares para orientar o planejamento.
Ela não representa uma avaliação completa dos riscos da entidade, pois
não houve acesso aos controles internos, aos sistemas contábeis nem aos
responsáveis pela preparação das demonstrações.

\section{6. Procedimentos de auditoria e testes
analíticos}\label{procedimentos-de-auditoria-e-testes-analuxedticos}

\subsection{6.1. Visão geral}\label{visuxe3o-geral}

Os procedimentos desenvolvidos foram organizados para testar a
integridade e a consistência da base sintética criada para o trabalho. A
finalidade é demonstrar a aplicação de técnicas analíticas e a
documentação de resultados de forma reproduzível.

A execução sobre dados artificiais não equivale à inspeção de documentos
originais nem à obtenção de evidência sobre transações reais da
companhia. Por isso, os resultados são apresentados como resultados de
testes sintéticos.

\subsection{6.2. Integridade da
população}\label{integridade-da-populauxe7uxe3o}

A população sintética contém 500.000 registros. Os testes de integridade
verificam a quantidade mínima prevista, a unicidade dos identificadores
e a identificação dos registros como sintéticos.

Os controles executados indicaram:

\begin{itemize}
\tightlist
\item
  500.000 registros na base;
\item
  500.000 identificadores únicos;
\item
  500.000 registros classificados como sintéticos;
\item
  conciliação do saldo registrado agregado com o saldo líquido público
  utilizado como referência.
\end{itemize}

Esses controles indicam que os requisitos estruturais definidos para a
população sintética foram atendidos. Eles não comprovam a integridade de
uma população contábil real da companhia.

\subsection{6.3. Teste de recálculo}\label{teste-de-recuxe1lculo}

O teste de recálculo compara, para cada registro sintético, o valor
registrado com o valor recalculado de acordo com as regras implementadas
no código.

O objetivo é identificar diferenças aritméticas e avaliar se o
procedimento consegue localizar inconsistências deliberadamente
inseridas na base.

A execução identificou 250 registros com diferenças entre os valores
registrados e os valores recalculados. O saldo registrado agregado foi
de R\$ 15.199.199 mil, enquanto o saldo recalculado agregado foi de R\$
15.216.626,8 mil. Considerando a diferença como saldo registrado menos
saldo recalculado, o resultado foi de −R\$ 17.427,8 mil, equivalente a
uma diferença absoluta de R\$ 17.427,8 mil.

Os valores representam resultados da simulação. Não constituem
distorções identificadas nas demonstrações contábeis ou nos registros
reais da Equatorial Energia Distribuição S.A.

\subsection{6.4. Verificação da documentação de
suporte}\label{verificauxe7uxe3o-da-documentauxe7uxe3o-de-suporte}

Foi executado um teste para identificar registros sintéticos sem
indicação de documento de suporte. O resultado apontou 5.459 registros
nessa condição.

A ausência de um identificador de suporte na base sintética é tratada
como uma exceção para fins do exercício. Em uma auditoria real, seria
necessário verificar a natureza do registro, a documentação exigível e a
existência de evidência alternativa antes de concluir que há uma
deficiência documental.

Portanto, o resultado não demonstra ausência de documentos reais da
companhia.

\subsection{6.5. Teste do período de
referência}\label{teste-do-peruxedodo-de-referuxeancia}

O teste de período identificou 75 registros sintéticos com datas fora do
exercício de 2025.

O objetivo é demonstrar um procedimento de verificação do corte temporal
e identificar registros que mereceriam investigação adicional caso a
população fosse real.

A presença de datas fora do período de referência não caracteriza
automaticamente erro contábil: seria necessário verificar a natureza da
transação, o período a que ela pertence e os critérios de reconhecimento
aplicáveis.

\subsection{6.6. Amostragem ilustrativa}\label{amostragem-ilustrativa}

Foi gerada uma amostra ilustrativa de 100 registros para demonstrar a
organização de um procedimento de amostragem.

A seleção não deve ser interpretada como amostra representativa dos
registros reais da companhia. Além disso, a adequação de um tamanho de
amostra depende do objetivo do teste, da população, do risco avaliado,
da distorção tolerável e de outros fatores de planejamento.

Em uma auditoria real, o auditor deveria documentar o método de seleção,
a definição da população, a unidade amostral, o objetivo do procedimento
e a avaliação dos resultados, conforme os requisitos aplicáveis.

\subsection{7. Resultados dos
procedimentos}\label{resultados-dos-procedimentos}

Os procedimentos analíticos foram executados sobre uma população
sintética de 500.000 registros, gerada em R para fins didáticos. Os
resultados a seguir referem-se exclusivamente a essa base e não
constituem achados reais da Equatorial Energia Distribuição S.A.

\subsubsection{7.1. Integridade da população e
conciliação}\label{integridade-da-populauxe7uxe3o-e-conciliauxe7uxe3o}

A população apresentou 500.000 registros e nenhum identificador
duplicado. Os saldos registrados na base sintética conciliam com os
valores de referência utilizados na simulação:

{\def\LTcaptype{none} % do not increment counter
\begin{longtable}[]{@{}lr@{}}
\toprule\noalign{}
Verificação & Resultado \\
\midrule\noalign{}
\endhead
\bottomrule\noalign{}
\endlastfoot
Ativo financeiro bruto & R\$ 21.062.007 mil \\
Obrigações especiais vinculadas & −R\$ 5.862.808 mil \\
Saldo líquido de referência & R\$ 15.199.199 mil \\
Diferença entre saldo líquido registrado e referência & R\$ 0 \\
\end{longtable}
}

A conciliação demonstra que os totais da base foram construídos para
corresponder aos valores de referência publicados. Não comprova a
exatidão dos registros individuais reais nem substitui a obtenção de
evidência de auditoria apropriada e suficiente.

\subsubsection{7.2. Diferenças individuais de
recálculo}\label{diferenuxe7as-individuais-de-recuxe1lculo}

Foram inseridas intencionalmente 250 diferenças entre os saldos
registrados e recalculados dos registros sintéticos. O teste identificou
250 registros com diferenças:

\begin{itemize}
\tightlist
\item
  Erros sintéticos inseridos: 250.
\item
  Erros sintéticos detectados: 250.
\item
  Erros inseridos não detectados: 0.
\item
  Diferenças detectadas que não constavam no controle de erros
  inseridos: 0.
\end{itemize}

O saldo registrado agregado foi de R\$ 15.199.199 mil, enquanto o saldo
recalculado agregado foi de R\$ 15.216.626,8 mil. Considerando a
diferença como saldo registrado menos saldo recalculado, o resultado foi
de −R\$ 17.427,8 mil, equivalente a uma diferença absoluta de R\$
17.427,8 mil. Trata-se de um resultado da simulação; diferenças
individuais positivas e negativas podem compensar-se parcialmente no
total.

O fato de todos os erros deliberadamente inseridos terem sido detectados
demonstra o funcionamento do procedimento nessa base específica. Não
permite concluir que o procedimento detectaria todas as distorções
possíveis em dados reais.

\subsubsection{7.3. Exceções documentais e
temporais}\label{exceuxe7uxf5es-documentais-e-temporais}

Também foram identificadas as seguintes situações sintéticas:

\begin{itemize}
\tightlist
\item
  5.459 registros classificados como sem documento de suporte.
\item
  75 registros com datas fora do exercício de 2025.
\end{itemize}

Essas situações foram construídas artificialmente para fins didáticos.
Em um trabalho real, seria necessário obter documentação original,
examinar as causas e avaliar eventuais impactos sobre os saldos e
divulgações. Como esses elementos não foram disponibilizados, não é
possível concluir que existam exceções equivalentes nos registros reais
da companhia.

\subsubsection{7.4. Síntese dos
resultados}\label{suxedntese-dos-resultados}

Os testes demonstram que os scripts conseguem verificar a integridade
básica da população construída, conciliar seus totais com os valores de
referência e sinalizar as inconsistências inseridas artificialmente. O
resultado constitui uma demonstração acadêmica de procedimentos
analíticos, não uma conclusão de auditoria sobre os ativos financeiros
de concessão ou as obrigações especiais reais da entidade.

\subsection{8. Papéis de trabalho e
documentação}\label{papuxe9is-de-trabalho-e-documentauxe7uxe3o}

O papel de trabalho PT 500-530-01 documenta os procedimentos de
integridade, conciliação e testes analíticos realizados sobre a
população sintética de 500.000 registros. A documentação está organizada
em três frentes distintas, para evitar que a conciliação dos saldos de
referência seja confundida com as exceções inseridas artificialmente.

\subsubsection{8.1. Conciliação dos saldos de
referência}\label{conciliauxe7uxe3o-dos-saldos-de-referuxeancia}

Os saldos registrados na base sintética foram comparados com os valores
de referência das demonstrações financeiras de 2025:

\begin{itemize}
\tightlist
\item
  Ativo financeiro bruto: R\$ 21.062.007 mil.
\item
  Obrigações especiais vinculadas: −R\$ 5.862.808 mil.
\item
  Ativo financeiro líquido de referência: R\$ 15.199.199 mil.
\end{itemize}

A população sintética concilia com esses valores de referência. Essa
conciliação demonstra a consistência dos totais construídos para a
simulação, mas não comprova a exatidão dos registros contábeis reais da
companhia.

\subsubsection{8.2. Diferenças individuais de
recálculo}\label{diferenuxe7as-individuais-de-recuxe1lculo-1}

Foram inseridas artificialmente 250 diferenças entre os saldos
registrados e os valores recalculados de registros sintéticos. O
procedimento analítico identificou 250 registros com diferenças.

A diferença agregada entre o saldo registrado e o saldo recalculado foi
de R\$ 17.427,8 mil. Esse valor representa exclusivamente o resultado da
simulação. Como as diferenças individuais podem ter sinais positivos ou
negativos, o valor agregado não corresponde à soma dos valores absolutos
das exceções.

A comparação com o controle de erros inseridos demonstrou que os 250
erros artificiais foram detectados, sem erros inseridos não detectados e
sem diferenças adicionais de recálculo fora do controle. Esse resultado
avalia o funcionamento do teste na base construída, mas não permite
estimar a eficácia do procedimento em registros reais.

\subsubsection{8.3. Exceções documentais e
temporais}\label{exceuxe7uxf5es-documentais-e-temporais-1}

A base sintética contém 5.459 registros classificados como sem documento
de suporte e 75 registros com datas fora do exercício de 2025. Essas
ocorrências foram criadas para fins didáticos.

Em uma auditoria real, situações semelhantes exigiriam obtenção e exame
da documentação original, avaliação da competência e investigação das
causas. Como não houve acesso a documentos internos ou sistemas da
companhia, não é possível concluir que existam exceções equivalentes nos
registros reais da Equatorial Energia Distribuição S.A.

\subsubsection{8.4. Referências e
limitações}\label{referuxeancias-e-limitauxe7uxf5es}

Os arquivos \texttt{resultados/07\_PT500\_PT530.csv} e
\texttt{resultados/14\_conciliacao\_e\_excecoes\_separadas.csv}
apresentam a documentação detalhada da conciliação, das diferenças
individuais e das demais exceções simuladas.

Os procedimentos foram executados sobre dados sintéticos gerados em R.
Não foram examinados o razão contábil real, contratos, documentos
originais, sistemas internos ou confirmações externas da companhia.
Portanto, os resultados não constituem evidência de auditoria real nem
fundamento para emitir opinião sobre as demonstrações financeiras.

\section{9. Programas de auditoria por
área}\label{programas-de-auditoria-por-uxe1rea}

\subsection{9.1. Objetivo dos programas}\label{objetivo-dos-programas}

Os programas de auditoria estabelecem procedimentos que poderiam ser
aplicados às áreas relevantes das demonstrações contábeis. Sua
finalidade é relacionar os riscos identificados aos testes planejados,
às evidências necessárias e às conclusões que deverão ser documentadas.

Neste trabalho, os programas foram elaborados como planejamento
acadêmico. Exceto pelos procedimentos expressamente identificados como
executados sobre a base sintética, os testes descritos nesta seção não
foram realizados em registros reais da companhia.

\subsection{9.2. Receita operacional}\label{receita-operacional}

\textbf{Objetivo:} avaliar, em um cenário de auditoria real, se as
receitas foram reconhecidas de acordo com os critérios contábeis
aplicáveis, no período correto e com documentação adequada.

\textbf{Procedimentos planejados:} - Compreender as políticas de
reconhecimento de receita e os processos relevantes. - Relacionar os
riscos de fraude e de distorção relevante às afirmações contábeis
pertinentes. - Selecionar transações e confrontar os registros com
contratos, documentos de faturamento e evidências de prestação do
serviço. - Examinar transações próximas ao encerramento do exercício
para avaliar o corte. - Desenvolver procedimentos analíticos e
investigar variações incomuns. - Avaliar a apresentação e a divulgação
das receitas nas demonstrações contábeis.

\textbf{Situação:} planejado; não executado em registros reais.

\subsection{9.3. Caixa e equivalentes de
caixa}\label{caixa-e-equivalentes-de-caixa}

\textbf{Objetivo:} avaliar a existência, a integralidade, a
titularidade, a exatidão e a apresentação dos saldos de caixa e
equivalentes.

\textbf{Procedimentos planejados:} - Obter a composição dos saldos e
conciliar os registros contábeis com extratos bancários. - Examinar
conciliações bancárias e investigar itens pendentes relevantes. -
Planejar confirmações externas com instituições financeiras, quando
aplicável. - Verificar restrições de uso, garantias e classificações
relevantes. - Avaliar a adequada apresentação e divulgação dos saldos.

\textbf{Situação:} planejado; não executado em registros reais. A
inclusão deste programa no dossiê não significa que a relevância dos
estoques tenha sido confirmada para o escopo analisado. Essa avaliação
dependeria do acesso às demonstrações contábeis, às informações
detalhadas da entidade e à definição do escopo de auditoria.

\subsection{9.4. Estoques}\label{estoques}

\textbf{Objetivo:} avaliar a existência, a integralidade, a mensuração e
a apresentação dos estoques, quando relevantes para as demonstrações
contábeis examinadas.

\textbf{Procedimentos planejados:} - Obter a composição dos estoques e
reconciliá-la com a contabilidade. - Avaliar os critérios de mensuração
e a consistência das políticas contábeis. - Planejar observação de
inventário físico e testes de contagem, conforme aplicável. - Examinar
itens obsoletos, danificados ou de baixa movimentação. - Testar
transações próximas ao encerramento do exercício. - Avaliar as
divulgações pertinentes.

\textbf{Situação:} planejado; não executado em registros reais. A
aplicabilidade e a relevância da área deverão ser confirmadas a partir
das informações da entidade.

\subsection{9.5. Imobilizado e ativos relacionados à
concessão}\label{imobilizado-e-ativos-relacionados-uxe0-concessuxe3o}

\textbf{Objetivo:} avaliar o reconhecimento, a mensuração, a
classificação, a movimentação e a divulgação dos ativos relacionados às
operações de concessão.

\textbf{Procedimentos planejados:} - Obter a composição e a movimentação
dos saldos contábeis. - Conciliar saldos iniciais, adições, baixas,
transferências, amortizações ou depreciações e saldos finais, conforme a
natureza do ativo. - Inspecionar contratos, documentos de aquisição,
memórias de cálculo e registros de movimentação. - Avaliar a
classificação entre ativo financeiro da concessão, ativo contratual e
ativo intangível, de acordo com os critérios aplicáveis. - Examinar as
obrigações especiais e sua relação com os ativos correspondentes. -
Avaliar premissas, estimativas e divulgações relevantes. - Investigar
variações significativas entre períodos.

\textbf{Situação:} foram executados testes analíticos exclusivamente
sobre a base sintética de 500.000 registros, incluindo verificações de
integridade, conciliação dos saldos agregados, recálculo e identificação
de exceções artificiais. Esses procedimentos demonstram a aplicação dos
scripts à população simulada, mas não constituem testes de auditoria
sobre os registros contábeis reais da companhia. Os procedimentos que
dependem de documentos originais, contratos, registros contábeis e
outras evidências reais permanecem planejados e não executados.

\subsection{9.6. Passivos e fornecedores}\label{passivos-e-fornecedores}

\textbf{Objetivo:} avaliar a integralidade, a existência, a exatidão, a
classificação e a apresentação dos passivos relevantes.

\textbf{Procedimentos planejados:} - Conciliar os saldos auxiliares com
o razão contábil. - Examinar documentos de suporte, contratos, faturas e
comprovantes pertinentes. - Investigar pagamentos posteriores ao
encerramento para identificar possíveis obrigações não registradas. -
Planejar confirmações externas quando apropriado. - Examinar o corte de
despesas e obrigações. - Avaliar classificações, vencimentos,
contingências e divulgações pertinentes.

\textbf{Situação:} planejado; não executado em registros reais.

\subsection{9.7. Estimativas contábeis}\label{estimativas-contuxe1beis}

\textbf{Objetivo:} avaliar a razoabilidade das estimativas contábeis e a
adequação das respectivas divulgações.

\textbf{Procedimentos planejados:} - Identificar estimativas relevantes
e compreender os métodos utilizados. - Avaliar a adequação dos dados,
das premissas e dos métodos adotados pela administração. - Considerar
incertezas de mensuração e possíveis indicadores de viés. - Desenvolver,
quando apropriado, estimativa pontual ou intervalo independente para
comparação. - Examinar a coerência entre as estimativas, as informações
disponíveis e as divulgações.

\textbf{Situação:} planejado; não executado em estimativas ou registros
reais da companhia.

\section{10. Cobertura das Normas Brasileiras de
Contabilidade}\label{cobertura-das-normas-brasileiras-de-contabilidade}

\subsection{10.1. Abordagem adotada}\label{abordagem-adotada}

A cobertura normativa é organizada com base nas normas indicadas no
roteiro acadêmico. O objetivo é demonstrar como cada tema se relaciona
ao planejamento, à execução, à documentação ou à comunicação dos
trabalhos de auditoria.

A matriz de cobertura deve funcionar como instrumento de
rastreabilidade. Para cada norma, devem ser registrados os itens
pertinentes, o objetivo, o procedimento ou a atividade correspondente, a
evidência documental e a situação do atendimento.

A presença de uma norma no relatório não significa que seus requisitos
tenham sido integralmente cumpridos. A versão vigente, os itens
aplicáveis e eventuais revisões devem ser confirmados na fonte oficial
antes da conclusão do dossiê.

\subsection{10.2. Relação das normas com o
trabalho}\label{relauxe7uxe3o-das-normas-com-o-trabalho}

{\def\LTcaptype{none} % do not increment counter
\begin{longtable}[]{@{}
  >{\raggedright\arraybackslash}p{(\linewidth - 4\tabcolsep) * \real{0.3333}}
  >{\raggedright\arraybackslash}p{(\linewidth - 4\tabcolsep) * \real{0.3333}}
  >{\raggedright\arraybackslash}p{(\linewidth - 4\tabcolsep) * \real{0.3333}}@{}}
\toprule\noalign{}
\begin{minipage}[b]{\linewidth}\raggedright
Norma ou referencial
\end{minipage} & \begin{minipage}[b]{\linewidth}\raggedright
Tema principal no contexto do trabalho
\end{minipage} & \begin{minipage}[b]{\linewidth}\raggedright
Aplicação ou documentação prevista
\end{minipage} \\
\midrule\noalign{}
\endhead
\bottomrule\noalign{}
\endlastfoot
Estrutura Conceitual & Conceitos e elementos dos trabalhos de
asseguração & Caracterização do objeto, dos critérios, das
responsabilidades e do nível de asseguração \\
NBC TA 200 & Objetivos gerais do auditor independente & Discussão de
segurança razoável, ceticismo e julgamento profissional \\
NBC TA 210 & Termos do trabalho de auditoria & Elaboração acadêmica dos
termos e das precondições da auditoria \\
NBC TA 220 & Gestão da qualidade da auditoria & Planejamento de direção,
supervisão, revisão e recursos \\
NBC TA 230 & Documentação de auditoria & Organização dos papéis de
trabalho e referências cruzadas \\
NBC TA 240 & Responsabilidade do auditor relacionada à fraude &
Discussão de riscos de fraude e respostas planejadas \\
NBC TA 250 & Leis e regulamentos & Planejamento de identificação de
requisitos legais e regulatórios relevantes \\
NBC TA 260 & Comunicação com os responsáveis pela governança &
Preparação de comunicação acadêmica sobre alcance, riscos e
constatações \\
NBC TA 265 & Comunicação de deficiências de controle interno &
Estruturação de matriz de deficiências hipotéticas, sem apresentá-las
como reais \\
NBC TA 300 & Planejamento da auditoria & Organização da estratégia e dos
programas de auditoria \\
NBC TA 315 & Identificação e avaliação de riscos & Matriz preliminar de
riscos e afirmações contábeis \\
NBC TA 320 & Materialidade & Definição e documentação de parâmetros
ilustrativos \\
NBC TA 330 & Respostas aos riscos avaliados & Relação entre riscos e
procedimentos planejados \\
NBC TA 402 & Entidade que utiliza organização prestadora de serviços &
Avaliação de aplicabilidade e planejamento, se pertinente \\
NBC TA 450 & Avaliação de distorções identificadas & Documentação e
avaliação das diferenças encontradas na simulação \\
NBC TA 500 & Evidência de auditoria & Avaliação metodológica da natureza
e suficiência da evidência \\
NBC TA 501 & Considerações específicas sobre itens selecionados &
Avaliação de aplicabilidade e planejamento de procedimentos
específicos \\
NBC TA 505 & Confirmações externas & Planejamento de confirmações,
quando apropriado \\
NBC TA 510 & Trabalhos iniciais e saldos iniciais & Avaliação da
aplicabilidade ao contexto e ao escopo acadêmico \\
NBC TA 520 & Procedimentos analíticos & Execução de análises
quantitativas sobre a base sintética \\
NBC TA 530 & Amostragem em auditoria & Documentação da amostra
ilustrativa e de suas limitações \\
NBC TA 540 & Auditoria de estimativas contábeis & Planejamento de
avaliação de métodos, premissas e incertezas \\
NBC TA 550 & Partes relacionadas & Avaliação de aplicabilidade e
planejamento de procedimentos \\
NBC TA 560 & Eventos subsequentes & Planejamento de análise de eventos
posteriores ao encerramento \\
NBC TA 570 & Continuidade operacional & Discussão dos procedimentos
necessários em uma auditoria real \\
NBC TA 580 & Representações formais & Planejamento de solicitações
formais, sem simular declarações da administração \\
NBC TA 600 & Auditorias de grupos & Consideração do contexto consolidado
e da aplicabilidade da norma \\
\end{longtable}
}

\subsection{10.3. Situação da
cobertura}\label{situauxe7uxe3o-da-cobertura}

A matriz de cobertura do projeto deverá indicar separadamente:

\begin{itemize}
\tightlist
\item
  \textbf{Análise conceitual:} conteúdo estudado e relacionado ao
  planejamento.
\item
  \textbf{Procedimento executado em dados sintéticos:} teste realizado
  exclusivamente sobre a base artificial.
\item
  \textbf{Procedimento planejado:} atividade ainda não realizada,
  normalmente por depender de registros, documentos ou acesso não
  disponível.
\item
  \textbf{Pendente de verificação:} item que exige confirmação da versão
  normativa, da aplicabilidade ou da documentação correspondente.
\end{itemize}

A situação definitiva de cada norma deverá ser conferida com o arquivo
de controle de cobertura e os papéis de trabalho. Não se deve declarar
atendimento integral apenas porque existe uma descrição da norma no
relatório.

\section{11. Limitações do trabalho}\label{limitauxe7uxf5es-do-trabalho}

O presente estudo possui limitações que devem ser consideradas na
interpretação dos resultados.

Primeiramente, não houve acesso ao razão contábil original, aos
documentos de suporte das transações, aos contratos internos completos,
aos sistemas da companhia ou aos responsáveis pela preparação das
demonstrações contábeis. Também não foram realizadas confirmações
externas, inspeções físicas, entrevistas ou observações de controles
internos.

Em segundo lugar, a base analítica de 500.000 registros é sintética.
Embora tenha sido construída para conciliar seus totais aos saldos
públicos de referência e para permitir testes reproduzíveis, seus
registros individuais não representam transações reais da companhia.

Em terceiro lugar, as exceções identificadas decorrem das regras e das
inconsistências deliberadamente inseridas na simulação. Os resultados
não permitem concluir que a companhia possua erros, falhas de controle
ou documentação ausente.

Em quarto lugar, a materialidade utilizada é ilustrativa e foi definida
a partir de uma premissa acadêmica. Não corresponde necessariamente à
materialidade estabelecida pelos auditores independentes da entidade.

Por fim, a cobertura das normas e os procedimentos planejados precisam
ser conferidos quanto à versão vigente, à aplicabilidade e à evidência
documental. A conclusão deste relatório deverá respeitar essas
limitações e não poderá ser apresentada como opinião de auditoria
independente.

\section{12. Conclusão técnica}\label{conclusuxe3o-tuxe9cnica}

Com base nos procedimentos acadêmicos realizados, foi possível verificar
a integridade estrutural da população sintética de 500.000 registros,
sem identificadores duplicados, e a conciliação dos saldos registrados
na base com os valores de referência das demonstrações financeiras de
2025.

Nos testes analíticos sobre os dados sintéticos, foram identificadas 250
diferenças individuais de recálculo, correspondentes às 250
inconsistências inseridas deliberadamente para testar os procedimentos.
Todas foram detectadas pelo script, sem diferenças inseridas não
detectadas. A diferença agregada entre os saldos registrados e
recalculados foi de −R\$ 17.427,8 mil, considerando a diferença como
saldo registrado menos saldo recalculado. Esse resultado decorre
exclusivamente da base sintética.

Também foram identificados 5.459 registros sintéticos classificados como
sem documento de suporte e 75 registros com datas fora do exercício de
2025. Essas situações foram simuladas e não constituem achados reais da
companhia.

A conciliação dos saldos de referência, a detecção das inconsistências
inseridas e a identificação das demais exceções representam resultados
de testes sobre dados construídos para fins didáticos. Eles não
comprovam a exatidão dos registros contábeis reais, a eficácia dos
controles internos da entidade ou a inexistência de outras distorções.

Não houve acesso ao razão contábil, aos documentos originais, aos
contratos, aos sistemas internos ou a confirmações externas da
companhia. Por esse motivo, os procedimentos previstos para registros
reais permanecem planejados e não executados, e não é possível emitir
uma conclusão de auditoria sobre os ativos financeiros de concessão, as
obrigações especiais ou as demonstrações financeiras da Equatorial
Energia Distribuição S.A.

O trabalho deve ser entendido como um dossiê acadêmico de planejamento,
documentação normativa e demonstração de procedimentos analíticos sobre
dados sintéticos, sujeito à revisão das fontes públicas, das versões das
normas aplicáveis e dos demais documentos exigidos pela disciplina.

\section{13. Referências e documentos de
suporte}\label{referuxeancias-e-documentos-de-suporte}

\subsection{13.1. Demonstrações
contábeis}\label{demonstrauxe7uxf5es-contuxe1beis}

EQUATORIAL ENERGIA DISTRIBUIÇÃO S.A. \textbf{Demonstrações contábeis
individuais e consolidadas e notas explicativas referentes ao exercício
findo em 31 de dezembro de 2025.} Nota Explicativa nº 11 --- Ativo
financeiro da concessão. Valores expressos em milhares de reais. Página
impressa 44 e página 46 do arquivo PDF, conforme conferência da cópia
utilizada no trabalho.

Documento disponível em:
https://api.mziq.com/mzfilemanager/v2/d/62b21cba-838c-49a4-aaef-e0fb2350c169/2cf8b9a1-fd46-3e5f-51a4-058b301aaf01?origin=2

Data de acesso: 25 out. 2026.

\subsection{13.2. Normas de auditoria}\label{normas-de-auditoria}

CONSELHO FEDERAL DE CONTABILIDADE (CFC). \textbf{Normas Brasileiras de
Contabilidade --- Auditoria Independente.} Disponível no portal oficial
do CFC:
https://cfc.org.br/tecnica/normas-brasileiras-de-contabilidade/nbc-ta-de-auditoria-independente/.
A versão vigente de cada norma e a data de consulta devem ser
verificadas e registradas no controle de normas.

\subsection{13.3. Documentação produzida no
projeto}\label{documentauxe7uxe3o-produzida-no-projeto}

Integram o dossiê acadêmico, conforme aplicável:

\begin{itemize}
\tightlist
\item
  Base de dados sintética e arquivo compactado correspondente.
\item
  Dicionário de variáveis e descrição das premissas de geração.
\item
  Controle de inconsistências deliberadamente inseridas.
\item
  Scripts de geração, testes, análise de exceções, materialidade e
  riscos.
\item
  Arquivos de resultados dos procedimentos.
\item
  Papéis de trabalho e matriz de cobertura normativa.
\item
  Programas de auditoria por área.
\item
  Controle de fontes, controle de normas, índice do dossiê e instruções
  de reprodução.
\end{itemize}

Os arquivos de suporte deverão ser revisados para assegurar consistência
entre os resultados reproduzidos pelos scripts e os valores apresentados
neste relatório.

\begin{center}\rule{0.5\linewidth}{0.5pt}\end{center}

\textbf{Nota de encerramento:} este documento é um relatório acadêmico
preliminar. Não constitui relatório de auditor independente, não
expressa opinião sobre as demonstrações contábeis da Equatorial Energia
Distribuição S.A. e não deve ser interpretado como evidência de que os
procedimentos planejados foram executados em registros reais.

\end{document}
